\chapter{Anwendung des Modells und Handlungsempfehlungen}
\label{ch:handlungsempfehlungen}
\section*{Tagelöhner}
Neben keiner Arbeit scheint auch Arbeit als Tagelöhner ein signifikantes Risiko für Armut zu sein.
Unser Modell beschreibt, dass eine Gesellschaft, in der ungesicherte Arbeitsverhältnisse weit verbreitet sind, viel Armut und Ungleichheit aufweist.
Nachdem die starken Gewerkschaften des 20.\ Jahrhunderts in den letzten 40 Jahren stark an Bedeutung verloren haben, sehen
wir heute einen starken Shift zur sog.\ Gig-Economy, die eine moderne Ausprägung von Tagesarbeit darstellen könnte.
Sollten diese Trends andauern, lässt sich mit unserem Modell ein wachsendes Armutsproblem prognostizieren.

Allerdings ist wichtig anzumerken, dass diese Prognose auf einem Modell einer geschichtlichen Gesellschaft beruht.
Somit beschränkt sich der Vergleich an dieser Stelle auf die qualitative Ebene,
inwiefern ein Schluss von damals auf heute möglich ist, lässt sich nicht quantifizieren.

\section*{Frauen}
In unserem Modell sind Frauen als Familienoberhaupt stärker von Armut gefährdet als Familien mit Männern als Familienoberhaupt.
Zwar lässt sich hier eine parallele zu den durchschnittlich geringeren Einkommen von Frauen in der heutigen Gesellschaft vermuten,
unserer Ansicht nach ist die höhere Armutsrate an dieser Stelle jedoch insbesondere der Form der erhobenen Daten geschuldet.
So sind Frauen nur als Familienoberhaupt (also als Haupteinkommensquelle der Familie) vermerkt, wenn kein Mann älter als 20 Jahre in diesem Haushalt lebt.
Wir sehen somit nicht unbedingt den Unterschied zwischen Einkommen von Männern und Frauen, sondern hauptsächlich, dass Haushalte mit maximal einem Einkommen
eine höhere Armutsrate aufweisen als solche mit einem oder mehr Einkommen, was zu erwarten ist.
Wir können somit keine Aussage über den Stand von Frauen in der Arbeitswelt und der Gesellschaft im viktorianischen London treffen.

